\documentclass[10pt,a4paper]{article}
\usepackage[utf8]{inputenc}
\usepackage[T1]{fontenc}
\usepackage{lmodern}
\usepackage[left=18mm,right=18mm,top=14mm,bottom=14mm]{geometry}
\usepackage{xcolor}
\usepackage{hyperref}
\usepackage{tabularx}
\usepackage{enumitem}

\definecolor{cvblue}{RGB}{22,66,125}

\hypersetup{
    colorlinks=true,
    linkcolor=cvblue,
    urlcolor=cvblue
}

\pagestyle{empty}
\setlength{\parindent}{0pt}
\setlength{\parskip}{0pt}

% Section styling matching Broteen Das's CV
\newcommand{\cvsection}[1]{%
  \vspace{7pt}%
  {\color{cvblue}\itshape\large #1}\\[-4pt]%
  {\color{cvblue}\rule{\textwidth}{0.75pt}}\vspace{3pt}%
}

% Item header macros
\newcommand{\entryheader}[4]{%
  \textbf{#1} \hfill \textit{#2}\\%
  \textit{#3} \hfill \textit{#4}\\%
}

\begin{document}

% ─────────────────────────────────────────────────────────
% HEADER
% ─────────────────────────────────────────────────────────
\begin{center}
  {\Huge\itshape\color{cvblue} Abhinav Goel}\\[4pt]
  \textit{Residence/Domicile:} Muzaffarnagar, Uttar Pradesh, India - 251001\\
  \textit{E-mail:} \href{mailto:abhinav.goel.robotics@gmail.com}{abhinav.goel.robotics@gmail.com} \ \textcolor{cvblue}{\ding{96}} \ \textit{Phone number:} +91-8384857368\\
  \textit{Place of birth:} Muzaffarnagar, Uttar Pradesh, India \ \textcolor{cvblue}{\ding{96}} \ \textit{Date of birth:} 18-12-2004\\
  \textit{LinkedIn:} \href{https://www.linkedin.com/in/abhinav-g-1b878a2b5}{Abhinav Goel} \ \textcolor{cvblue}{\ding{96}} \ \textit{GitHub:} \href{https://github.com/Humobot1812}{Humobot1812} \ \textcolor{cvblue}{\ding{96}} \ \textit{Portfolio:} \href{https://Humobot1812.github.io/Portfolio/}{Humobot1812.github.io/Portfolio}
\end{center}

% ─────────────────────────────────────────────────────────
% EDUCATION
% ─────────────────────────────────────────────────────────
\cvsection{Education}

\textbf{B. Tech. in Electronics and Communication Engineering} \hfill \textit{PDPM IIITDMJ}\\
\textit{Bachelor's Degree Program} \hfill \textit{August 2024 - Present}\\
\textit{Semester IV (CPI - 8.5)}\\[4pt]

\textbf{Senior Secondary (Higher Secondary)} \hfill \textit{PR Public School, Muzaffarnagar}\\
\textit{Central Board of Secondary Education (CBSE)} \hfill \textit{July 2022 - July 2024}\\
\textit{Percentage: 93.6\%}\\[4pt]

\textbf{Secondary (High School)} \hfill \textit{Shardein School, Muzaffarnagar}\\
\textit{Central Board of Secondary Education (CBSE)} \hfill \textit{April 2020 - July 2022}\\
\textit{Percentage: 98.4\%}

% ─────────────────────────────────────────────────────────
% AREAS OF INTEREST
% ─────────────────────────────────────────────────────────
\cvsection{Areas of interest}
\begin{itemize}[leftmargin=1.2rem,itemsep=1.5pt,topsep=2pt,parsep=0pt]
  \item Autonomous Mobile Robots (AMRs), SLAM and Navigation (ROS 2, Nav2)
  \item Unmanned Aerial Vehicles (UAVs) \& Autonomous Aerial Flight Systems
  \item Robotic Manipulation, Visual Servoing \& Kinematics
  \item Computer Vision, Spatial Geolocation \& Multi-Spectral Perception
  \item Tactical Ground Control Stations (GCS) \& Human-Machine Interfaces (Qt6 / C++)
  \item Embedded Systems, Flight Controllers \& Discrete Power Electronics
  \item Subterranean, Field Exploration \& Bio-Inspired Legged Robotics
\end{itemize}

% ─────────────────────────────────────────────────────────
% HONORS AND KEY ACHIEVEMENTS
% ─────────────────────────────────────────────────────────
\cvsection{Honors and key achievements}
\begin{itemize}[leftmargin=1.2rem,itemsep=2.5pt,topsep=2pt,parsep=0pt]
  \item \textbf{e-Yantra Robotics Competition (e-YRC 2025–26), IIT Bombay:} Top 10 National Finalist among 400+ participating teams across India in the \textit{KrishiCobot} theme, designing an autonomous agricultural harvesting mobile manipulator with ROS 2, UR5 visual servoing, and precision vision pipelines. [\href{https://Humobot1812.github.io/Portfolio/assets/Certificates/eyantra_certificate.pdf}{Certificate}]
  \item \textbf{DD Robocon 2025 (Doordarshan \& IIT Delhi):} Achieved Stage-2 National Qualification representing Team Barbarick with an evaluation score of 86/100, designing competition-grade holonomic robot chassis and actuation mechanisms. [\href{https://Humobot1812.github.io/Portfolio/assets/Certificates/Robocon_certificate.jpeg}{Certificate}]
  \item \textbf{SAE Aerothon 2025 (SAE India):} Phase 2 National Qualification in the premier aerial robotics challenge, demonstrating an autonomous fixed-altitude surveillance quadcopter with ground target pixel-to-GPS conversion. [\href{https://Humobot1812.github.io/Portfolio/assets/Certificates/Aerothon_2025.png}{Certificate}]
  \item \textbf{NIDAR 2026 (National Innovation Challenge for Drone Application and Research):} National Finalist, Aero Fabrication Club, developing an autonomous disaster surveillance UAV with 75+ acre automated polygon coverage and real-time survivor detection. [\href{https://Humobot1812.github.io/Portfolio/assets/Certificates/Nidar_certificate.jpeg}{Certificate}]
\end{itemize}

% ─────────────────────────────────────────────────────────
% RESEARCH AND DEVELOPMENT EXPERIENCE
% ─────────────────────────────────────────────────────────
\cvsection{Research and development experience}

\textbf{Autonomous Agricultural Multi-Robot System (KrishiCobot)} \hfill \textit{August 2025 - May 2026}\\
\textit{e-Yantra Robotics Competition 2025–26, IIT Bombay} \hfill \textit{\href{https://github.com/Humobot1812/Krishi_Cobot}{github.com/Humobot1812/Krishi\_Cobot}}
\begin{itemize}[leftmargin=1.2rem,itemsep=1.5pt,topsep=2pt,parsep=0pt]
  \item Developed a multi-robot agricultural system integrating an autonomous mobile robot (eBot) for greenhouse navigation with a dock-based UR5 manipulator for automated fruit and fertilizer handling in Gazebo.
  \item Implemented a LiDAR-based RANSAC shape detector converting polar scans to Cartesian clusters, fitting lines and verifying vertex angles to classify triangular (fertilizer) and square (bad health) plant markers.
  \item Engineered a vision pipeline utilizing ArUco markers, HSV color segmentation, contour centroid extraction, and IQR depth filtering, broadcasting dynamic TF2 transforms for targets and docking stations.
  \item Programmed UR5 delta-twist servo control and PID positioning with intermediate waypoints to avoid kinematic singularities; qualified as a Top 10 National Finalist among 400+ collegiate teams across India.
\end{itemize}

\vspace{3pt}
\textbf{Autonomous Warehouse Stock Auditing AMR \& Active Vision System (WareOps)} \hfill \textit{December 2025 - August 2026}\\
\textit{Robotics \& Autonomous Systems, PDPM IIITDMJ} \hfill \textit{\href{https://github.com/WareOps}{github.com/WareOps}}
\begin{itemize}[leftmargin=1.2rem,itemsep=1.5pt,topsep=2pt,parsep=0pt]
  \item Architected a distributed ROS 2 Humble AMR system splitting workloads over Wi-Fi between an onboard Raspberry Pi 4 (navigation/control) and a remote operator workstation (vision/logging/teleop).
  \item Developed a custom C++ ros2\_control hardware interface plugin (diffdrive\_arduino) over USB serial with an Arduino Nano running ROSArduinoBridge firmware for motor PID control and quadrature encoder feedback.
  \item Configured SLAM Toolbox for 2D mapping and Nav2 with AMCL for autonomous rack-to-rack path planning and obstacle avoidance using an RPLidar A1M8 in physical warehouse environments.
  \item Engineered an active-vision scanning mast with an ESP32-CAM Wi-Fi MJPEG stream and HTTP pan-tilt servo tracking, running OpenCV WeChat QR and PyZbar to decode rack barcodes and log records to Excel and CSV ledgers.
\end{itemize}

\vspace{3pt}
\textbf{CAVE DIVER — Autonomous 6-Wheeled Subterranean Exploration Rover} \hfill \textit{September 2026 - Present}\\
\textit{Smart India Hackathon 2026, PDPM IIITDMJ} \hfill \textit{\href{https://github.com/Humobot1812/CAVE_DIVER_6_Wheeled_Rover}{github.com/Humobot1812/CAVE\_DIVER\_6\_Wheeled\_Rover}}
\begin{itemize}[leftmargin=1.2rem,itemsep=1.5pt,topsep=2pt,parsep=0pt]
  \item Engineered an autonomous 6-wheeled robotic rover platform in Ignition Gazebo (Fortress) with custom Armo\_bot URDF and cave SDF worlds for GPS-denied subterranean reconnaissance and hazard monitoring.
  \item Integrated dual-spectrum visual inspection, capturing synchronized optical RGB (/camera/image\_raw) and LWIR thermal (/thermal/image\_raw) video feeds with tactical HUD overlays for survivor and obstacle detection.
  \item Developed an atmospheric hazard telemetry suite continuously monitoring six environmental parameters: Temperature, Relative Humidity, CO\textsubscript{2}, O\textsubscript{2}, Methane (CH\textsubscript{4}), and Carbon Monoxide (CO).
  \item Built a web Ground Control Station in Flask 3.1 \& HTML5/WebSocket (port 5001) rendering live dual feeds, gas dials, a 2D LiDAR radar sweep display, spotlight control, and QR matrix landmark navigation via ros\_gz\_bridge.
\end{itemize}

\vspace{3pt}
\textbf{Tactical Multi-UAV Ground Control Station (DroneGCS)} \hfill \textit{September 2026 - Present}\\
\textit{Aero Fabrication Club, PDPM IIITDMJ} \hfill \textit{\href{https://github.com/Humobot1812/DroneGCS}{github.com/Humobot1812/DroneGCS}}
\begin{itemize}[leftmargin=1.2rem,itemsep=1.5pt,topsep=2pt,parsep=0pt]
  \item Developed a tactical Ground Control Station in C++17, Qt 6, and QML for multi-rotor, fixed-wing, and VTOL drones communicating via MAVLink v2 with ArduPilot and PX4 autopilots over UDP, TCP, and serial.
  \item Integrated a native Tactical AI Copilot powered by Google Gemini AI, enabling natural language flight commands, automated RTL execution, and conversational telemetry SITREP briefings.
  \item Built a Primary Flight Display (PFD) HUD with artificial horizon, Google Satellite imagery, and a live multi-channel Telemetry Oscilloscope plotting altitude, climb rate, battery voltage, and groundspeed.
  \item Engineered a Geofence Safety Console for setting cylindrical boundaries, altitude ceilings, and keep-in/out corridors, with automated synchronization from flight controllers and low-latency GStreamer H.264 video.
\end{itemize}

\vspace{3pt}
\textbf{SafePilot — Intelligent Flight Assistance \& Safety Envelope System} \hfill \textit{September 2026 - Present}\\
\textit{Open-Source UAV Development} \hfill \textit{\href{https://github.com/Humobot1812/SafePilot}{github.com/Humobot1812/SafePilot}}
\begin{itemize}[leftmargin=1.2rem,itemsep=1.5pt,topsep=2pt,parsep=0pt]
  \item Created an open-source ROS 2 Humble package (safe\_pilot) wrapping MAVSDK to provide beginner-safe flight assistance and software-enforced safety envelopes for ArduPilot and PX4 drones.
  \item Implemented 3D body-frame velocity teleoperation via gamepad/joystick with single-button discrete flight commands (Arm, Takeoff, Land, RTL, Hold, Kill) and live speed scaling via D-pad.
  \item Engineered a multi-core AI hand gesture control worker process with CPU affinity pinning and MediaPipe Hand Landmarker (TFLite XNNPACK delegate) for touchless strafe, climb, and closed-fist yaw rotation.
  \item Configured embedded polygon geofencing with altitude ceilings triggering automated auto-RTL on boundary breach, with one-command unified Gazebo Harmonic and ArduPilot SITL simulation launch.
\end{itemize}

\vspace{3pt}
\textbf{Autonomous Disaster Surveillance UAV \& Victim Detection} \hfill \textit{June 2025 - January 2026}\\
\textit{NIDAR 2026 National UAV Challenge, Aero Fabrication Club} \hfill \textit{\href{https://github.com/Humobot1812/Autonomous_drone}{github.com/Humobot1812/Autonomous\_drone}}
\begin{itemize}[leftmargin=1.2rem,itemsep=1.5pt,topsep=2pt,parsep=0pt]
  \item Developed an autonomous aerial search and rescue platform for the NIDAR 2026 National UAV Challenge, integrating computer vision, GPS mapping, and autonomous navigation using MAVSDK Python.
  \item Implemented an autonomous survey coverage system parsing .kml polygon boundary files to systematically sweep large disaster zones at 25 m altitude.
  \item Deployed real-time AI human detection models on live downward camera feeds, converting pixel detections to geographic GPS coordinates and dynamically controlling camera pan/tilt tracking toward victims.
  \item Validated complete mission execution (automated arming, takeoff, waypoint traversal, victim loiter, and RTL) across ArduPilot SITL simulation and physical multirotor hardware flight tests.
\end{itemize}

\vspace{3pt}
\textbf{Autonomous Aerial Survey \& Ground Target Geolocation Drone} \hfill \textit{May 2025 - November 2025}\\
\textit{SAE Aerothon 2025, Aero Fabrication Club, PDPM IIITDMJ} \hfill \textit{\href{https://github.com/Humobot1812/Aerothon-25}{github.com/Humobot1812/Aerothon-25}}
\begin{itemize}[leftmargin=1.2rem,itemsep=1.5pt,topsep=2pt,parsep=0pt]
  \item Engineered an autonomous survey and object detection drone using ROS 2, DroneKit, and Pixhawk for competitive presentation at SAE Aerothon 2025.
  \item Deployed YOLOv4-tiny using OpenCV DNN on an onboard companion computer for real-time edge target and hotspot identification during autonomous flight.
  \item Designed an automated mission interrupt pipeline: switches from AUTO to GUIDED mode upon target detection, descends, performs PID-based visual centering and hover inspection, then climbs and resumes survey.
  \item Derived an image pixel-to-GPS conversion algorithm considering camera FOV, altitude, and heading, applying Haversine spatial filtering ($<$0.8 m) to eliminate duplicate hotspot coordinate records.
\end{itemize}

\vspace{3pt}
\textbf{Differential Drive Mobile Robot — ROS 2 Navigation \& SLAM} \hfill \textit{June 2025 - July 2025}\\
\textit{Electronics and Robotics Society, PDPM IIITDMJ} \hfill \textit{\href{https://github.com/Humobot1812/Differential-Drive-ROS2-Navigation-Project}{Differential-Drive-ROS2}}
\begin{itemize}[leftmargin=1.2rem,itemsep=1.5pt,topsep=2pt,parsep=0pt]
  \item Constructed a complete simulation and navigation stack for a 4-wheeled differential-drive robot (Armo\_description) using ROS 2 Humble and Gazebo Harmonic.
  \item Modeled kinematic chains, wheel inertia matrices, collision geometries, 2D LiDAR, and IMU plugins using modular URDF/Xacro macros with gazebo\_bridge integration.
  \item Integrated robot\_localization Extended Kalman Filter (EKF) fusing wheel odometry and IMU data to produce drift-free robot state estimation in complex maze environments.
  \item Built 2D occupancy grid maps using SLAM Toolbox and configured Nav2 with AMCL particle filtering for autonomous obstacle avoidance, global path planning, and goal pose tracking.
\end{itemize}

\vspace{3pt}
\textbf{Smart Multi-Protocol ESP32 Robotics Transmitter System} \hfill \textit{April 2026 - July 2026}\\
\textit{Electronics and Robotics Society, PDPM IIITDMJ} \hfill \textit{\href{https://github.com/Humobot1812/Smart-Multi-Protocol-ESP32-Transmitter}{Smart-ESP32-Transmitter}}
\begin{itemize}[leftmargin=1.2rem,itemsep=1.5pt,topsep=2pt,parsep=0pt]
  \item Designed and built a modular smart transmitter utilizing a dual-ESP32 architecture: a Controller ESP32 reading inputs and a Dashboard ESP32 managing UI and RF transmission via UART.
  \item Interfaced 2 analog joysticks, 5 potentiometers, a 4×4 matrix keypad, and digital push buttons for comprehensive multichannel robotics teleoperation.
  \item Programmed a 128×64 ST7920 graphical LCD using the U8g2 library, featuring boot animations, keypad password security, real-time channel bar graphs, battery monitoring, and wheel ratio configuration.
  \item Supported four switchable wireless protocols (NRF24L01+ 2.4 GHz, ESP-NOW 250 m, Wi-Fi UDP, and Bluetooth), with an onboard MPU6050 IMU enabling tilt-based gesture control.
\end{itemize}

\vspace{3pt}
\textbf{Custom ESP32-S3 Quadcopter Drone with Discrete MOSFET ESCs} \hfill \textit{January 2026 - February 2026}\\
\textit{Aero Fabrication Club, PDPM IIITDMJ} \hfill \textit{\href{https://youtube.com/playlist?list=PLQcgql__dXrcVR0A4IEijRmWpTvdtAOpG}{ESP32-Quadcopter-Playlist}}
\begin{itemize}[leftmargin=1.2rem,itemsep=1.5pt,topsep=2pt,parsep=0pt]
  \item Designed and fabricated custom 3-phase Electronic Speed Controllers (ESCs) using discrete N-channel MOSFETs and gate drivers, eliminating commercial off-the-shelf ESC modules.
  \item Implemented 6-step trapezoidal commutation firmware on an ESP32-S3 microcontroller, achieving smooth BLDC motor speed regulation with configurable acceleration profiles.
  \item Programmed a 500 Hz flight stabilization PID loop sampling angular velocities and accelerations from an MPU6050 6-axis IMU for real-time roll, pitch, and yaw stabilization.
  \item Integrated a BMP280 barometric pressure sensor for relative altitude estimation and hover stabilization, validating performance through physical motor bench testing and flight trials.
\end{itemize}

\vspace{3pt}
\textbf{5-DOF Intelligent Robotic Arm with AI Hand Tracking \& ESP32-CAM} \hfill \textit{March 2025 - Present}\\
\textit{Electronics and Robotics Society, PDPM IIITDMJ} \hfill \textit{\href{https://github.com/Humobot1812/Robotic_arm}{github.com/Humobot1812/Robotic\_arm}}
\begin{itemize}[leftmargin=1.2rem,itemsep=1.5pt,topsep=2pt,parsep=0pt]
  \item Fabricated an articulated 5-DOF robotic arm (Base, Shoulder, Elbow, Wrist, Gripper) driven by an Arduino Uno interfacing with a PCA9685 16-channel I²C servo controller.
  \item Mounted an ESP32-CAM on the end-effector providing real-time wireless Wi-Fi video streaming of manipulation tasks to a custom desktop operator dashboard.
  \item Developed an AI vision teleoperation pipeline in Python using OpenCV and MediaPipe, tracking 21 hand landmarks to convert human hand motion into real-time arm servo commands.
  \item Implemented pinch gesture recognition for gripper open/close and geometric Forward and Inverse Kinematics (FK/IK) with workspace boundary checks and KP gain tuning over USB serial.
\end{itemize}

\vspace{3pt}
\textbf{Bio-Inspired Serpentine Snake Robot Simulation} \hfill \textit{June 2025 - July 2025}\\
\textit{Electronics and Robotics Society, PDPM IIITDMJ} \hfill \textit{\href{https://github.com/Humobot1812/Snake_Robot}{github.com/Humobot1812/Snake\_Robot}}
\begin{itemize}[leftmargin=1.2rem,itemsep=1.5pt,topsep=2pt,parsep=0pt]
  \item Engineered a bio-inspired multi-segment snake robot simulation in ROS 2 Humble and Gazebo utilizing modular snake\_robot\_description and snake\_robot\_controller packages.
  \item Constructed modular body segments interconnected with revolute joints in URDF/Xacro, configuring Gazebo physics plugins and ros\_gz\_bridge topic communications.
  \item Developed a custom Python controller node generating traveling sinusoidal undulation waves by publishing coordinated phase-shifted joint position commands to actuate body links.
  \item Implemented configurable motion parameters allowing real-time adjustment of wave amplitude, frequency, and inter-joint phase offsets for smooth snake-like planar locomotion.
\end{itemize}

\vspace{3pt}
\textbf{IoT Hostel Attendance System with ESP8266 \& RFID} \hfill \textit{October 2025 - November 2025}\\
\textit{Campus Automation Project, PDPM IIITDMJ} \hfill \textit{\href{https://github.com/Humobot1812/Hostel_Attendance_system}{Hostel\_Attendance\_system}}
\begin{itemize}[leftmargin=1.2rem,itemsep=1.5pt,topsep=2pt,parsep=0pt]
  \item Engineered an IoT-based hostel attendance system utilizing an ESP8266 (NodeMCU) and an MFRC522 RFID reader interfaced via SPI with MIFARE 1K RFID cards.
  \item Programmed automatic entry/exit state machine logic: initial RFID scan marks Time In and subsequent scan marks Time Out, tracking Student ID and gate numbers automatically.
  \item Integrated direct cloud logging: ESP8266 transmits HTTPS POST requests over Wi-Fi to a Google Apps Script Web App (sheet.gs), updating records in real-time Google Sheets without dedicated local servers.
  \item Developed companion firmware (rfid\_tag.ino) for writing student credentials onto RFID tags; deployed and validated system functionality in institute hostel operations.
\end{itemize}

% ─────────────────────────────────────────────────────────
% CERTIFICATIONS
% ─────────────────────────────────────────────────────────
\cvsection{Certifications}
\begin{itemize}[leftmargin=1.2rem,itemsep=2pt,topsep=2pt,parsep=0pt]
  \item \textbf{e-Yantra Robotics Competition (e-YRC 2025–26) Finalist:} Awarded by IIT Bombay for securing Top 10 nationally among 400+ collegiate teams in the autonomous KrishiCobot agricultural robotics theme (August 2025 – May 2026).
  \item \textbf{DD Robocon 2025 Stage-2 Certificate:} Awarded by Doordarshan \& IIT Delhi for qualifying Stage 2 with Team Barbarick, scoring 86/100 in technical design and mechanism evaluations (2025).
  \item \textbf{SAE Aerothon 2025 National Aerial Robotics Certificate:} Awarded by SAE India for qualifying Phase 2 in the national autonomous UAV surveillance challenge (June 2025).
  \item \textbf{NIDAR 2026 National UAV Challenge:} Certified participant in National Innovation Challenge for Drone Application and Research (January 2026).
  \item \textbf{Hands-on Introduction to Autonomous Systems with ROS 2 \& Embedded Hardware:} PDPM IIITDMJ Student Technical Workshop (2024).
\end{itemize}

% ─────────────────────────────────────────────────────────
% TECHNICAL SKILLS
% ─────────────────────────────────────────────────────────
\cvsection{Technical Skills}
\begin{tabularx}{\textwidth}{@{}p{5.6cm}X@{}}
  \textbf{Robotics \& Autonomous Systems:} & ROS 2 (Humble), Nav2, ros2\_control, SLAM Toolbox, AMCL, TF2, URDF/Xacro, MoveIt 2, Kinematics, Trajectory Planning, Sensor Fusion, EKF, RANSAC\\[3pt]
  \textbf{UAV \& Flight Systems:} & ArduPilot, PX4, Pixhawk, MAVLink, MAVSDK, DroneKit, QGroundControl, Mission Planner, ArduPilot SITL, PID Control\\[3pt]
  \textbf{Computer Vision \& Perception:} & OpenCV, YOLOv8, YOLOv4-Tiny, ArUco, Visual Servoing, Pinhole Camera Model, ArUco Marker Geolocation, ORB-SLAM3\\[3pt]
  \textbf{Embedded Systems \& Electronics:} & ESP32-S3, ESP32-CAM, Arduino, STM32, UART, I$^2$C, SPI, PWM, MPU9250/MPU6050, Encoder Motors, Barometer, NRF24L01, LoRa, ST7920 Display, Custom MOSFET ESC\\[3pt]
  \textbf{Compute, Sensors \& Hardware:} & Raspberry Pi 4/5, NVIDIA Jetson Nano/AGX Xavier, RPLidar A1M8, LiDAR, Raspberry Pi Camera v3 NoIR\\[3pt]
  \textbf{Simulation \& Engineering Tools:} & Gazebo, RViz2, MuJoCo, MATLAB/Simulink, Fusion 360, 3D Printing, TinkerCAD, Git, Linux (Ubuntu 22.04)\\[3pt]
  \textbf{Programming Languages:} & Python, C++\\
\end{tabularx}

% ─────────────────────────────────────────────────────────
% LANGUAGE PROFICIENCY
% ─────────────────────────────────────────────────────────
\cvsection{Language proficiency}
\begin{tabularx}{\textwidth}{@{}p{5.1cm}X@{}}
  \textbf{English, Hindi} & Full Professional Proficiency\\
\end{tabularx}

% ─────────────────────────────────────────────────────────
% MEMBERSHIPS
% ─────────────────────────────────────────────────────────
\cvsection{Memberships}
\begin{itemize}[leftmargin=1.2rem,itemsep=2pt,topsep=2pt,parsep=0pt]
  \item \textbf{Electronics and Robotics Society of PDPM IIITDMJ} (Coordinator, March 2026 – Present | Member, August 2024 – March 2026)
  \item \textbf{Aero Fabrication Club of PDPM IIITDMJ} (Core Committee Member, May 2025 – Present | Member, August 2024 – May 2025)
\end{itemize}

\end{document}
